\chapter{Introducción}

\section{Contexto y motivación}

La contaminación por residuos en entornos terrestres y acuáticos constituye uno de
los principales problemas ambientales actuales. La presencia de basura en ríos,
playas, zonas costeras, bosques, áreas urbanas y otros espacios naturales afecta
a los ecosistemas, a la biodiversidad y a la calidad de vida de las personas.
Para comprender la magnitud del problema y evaluar la eficacia de las medidas de
prevención o mitigación, las administraciones públicas, las instituciones
científicas y las organizaciones ambientales impulsan programas de monitorización
de residuos en distintos compartimentos ambientales.

La utilidad de estos programas depende de que los datos recogidos sean
comparables, estén correctamente georreferenciados y conserven información
suficiente sobre cómo, cuándo y dónde se han obtenido. En el caso de las basuras
fluviales, el Joint Research Centre de la Comisión Europea señala la falta de
información sistemática sobre la cantidad de residuos transportados por los ríos
y la ausencia de metodologías armonizadas capaces de generar datos cuantitativos
comparables \cite{gonzalez2016riverine}. De forma similar, trabajos recientes
subrayan la necesidad de armonizar las técnicas de monitorización de
macroplásticos en ecosistemas acuáticos y terrestres
\cite{gallitelli2024macroplastics}.

Un antecedente directo de este proyecto es el trabajo realizado para la
monitorización de macrobasuras flotantes en ríos, donde se desarrolló una
aplicación para dispositivos tipo \textit{tablet} capaz de registrar los ítems
observados, su tamaño y su posición geográfica durante sesiones de
monitorización \cite{gonzalezfernandez2017harmonized}. Este antecedente refuerza
la idea de que la herramienta de captura no debe limitarse a digitalizar un
formulario, sino que debe orientar la observación y conservar desde el inicio los
datos necesarios para que los registros puedan compararse posteriormente.

En este contexto, el cliente planteó la necesidad de renovar y ampliar la
aplicación \textit{Floating Litter Monitoring}. La herramienta original estaba
orientada principalmente a la monitorización de residuos en ríos y mares,
mientras que el nuevo proyecto debía extender esa capacidad a una mayor variedad
de entornos naturales y urbanos, incorporar una gestión más flexible de usuarios
y grupos, mejorar la recogida de evidencias y permitir la consulta posterior de
los datos desde un portal web.

Como respuesta a esta necesidad se propone el desarrollo de OmniLitter, una
plataforma formada por una aplicación móvil, un portal web y un backend
centralizado. Su finalidad es facilitar la recogida, gestión y análisis de datos
de residuos mediante una solución que combine trabajo de campo, sincronización de
información, control de permisos y explotación posterior de los datos
recopilados.

\section{Problema abordado}

La monitorización de residuos requiere la participación de distintas personas y
grupos que realizan campañas en lugares y momentos diferentes. Si la información
se registra mediante formularios no normalizados, hojas de cálculo aisladas o
campañas de limpieza sin un protocolo común, los datos resultantes suelen quedar
dispersos, son difíciles de comparar y pierden parte de su valor para análisis
científicos posteriores.

Además, la captura de datos se realiza habitualmente en la naturaleza, donde la
conectividad a Internet puede ser limitada o inexistente. Esta situación dificulta
el uso de soluciones que dependen de una conexión permanente con servidores
remotos. Por ello, resulta necesario que la aplicación móvil pueda almacenar
sesiones, observaciones, fotografías y metadatos en el propio dispositivo, y que
sincronice esa información con el servidor cuando vuelva a existir conexión.

El proyecto también exige registrar datos con suficiente calidad científica. Cada
observación debe asociarse a una taxonomía de residuos, una localización
geográfica, una fecha y hora, evidencias fotográficas y metadatos de muestreo
como el entorno, la tipología, el método empleado, la duración, el transecto o el
área observada. La precisión de la posición GPS y el estado de sincronización de
cada registro también son relevantes para interpretar correctamente los datos
recogidos.

A estas necesidades se suma la organización colaborativa del trabajo. Los
usuarios pertenecen a grupos, los grupos gestionan campañas, las campañas
contienen sesiones de muestreo y las sesiones agrupan observaciones. Esta
estructura requiere mecanismos de autenticación, control de permisos y aislamiento
de datos para que cada equipo pueda gestionar su información de forma segura y
privada.

Por tanto, el problema abordado no consiste únicamente en crear un formulario
digital de recogida de residuos. El reto es diseñar una plataforma que conecte la
captura offline en campo con una infraestructura común de almacenamiento,
administración, visualización y exportación, manteniendo la trazabilidad y la
comparabilidad de los datos desde su origen.

\section{Objetivos}

El objetivo principal de este Trabajo Fin de Grado es diseñar, desarrollar e
implementar OmniLitter, una plataforma software multiplataforma que facilite la
monitorización de residuos en distintos entornos naturales y permita transformar
observaciones de campo en datos estructurados, comparables y reutilizables.

Para alcanzar este objetivo general se establecen los siguientes objetivos
específicos:

\begin{enumerate}
    \item Diseñar una arquitectura software que integre una aplicación móvil, un
    portal web, un backend centralizado y una base de datos común.

    \item Organizar la información del sistema mediante grupos de trabajo,
    campañas, sesiones de muestreo y observaciones geolocalizadas.

    \item Desarrollar una aplicación móvil destinada a la captura de datos en
    campo, incluyendo residuos observados, fotografías, ubicación, fecha, hora y
    metadatos de muestreo.

    \item Permitir el funcionamiento de la aplicación móvil en condiciones de
    conectividad limitada mediante persistencia local y sincronización diferida
    con el servidor.

    \item Incorporar taxonomías estandarizadas, entornos, tipologías y métodos
    de muestreo configurables para favorecer la calidad científica de la
    información registrada.

    \item Implementar mecanismos de autenticación, roles, permisos y aislamiento
    de datos entre grupos de trabajo.

    \item Desarrollar un portal web que facilite la administración de usuarios,
    grupos y campañas, así como la consulta de sesiones y observaciones
    recopiladas.

    \item Proporcionar herramientas de visualización y explotación de datos,
    incluyendo mapas, paneles informativos y exportación en formatos adecuados
    para análisis posterior.

    \item Garantizar la integridad, trazabilidad y consistencia de la información
    almacenada mediante un backend centralizado y una API autenticada.

    \item Aplicar metodologías y herramientas de ingeniería del software para
    documentar, probar y validar la solución desarrollada dentro del alcance
    académico del proyecto.
\end{enumerate}

\section{Alcance del proyecto}

El presente Trabajo Fin de Grado aborda el diseño, desarrollo e implementación de
una versión funcional de una plataforma software denominada OmniLitter, orientada
a facilitar la monitorización y el análisis de la contaminación por residuos en
distintos ecosistemas terrestres, acuáticos y urbanos, dentro del alcance
académico del proyecto.

La solución desarrollada está compuesta por una aplicación móvil destinada a la
recogida de información en campo, un portal web para la gestión y explotación de
los datos recopilados, y un backend centralizado encargado de la autenticación,
la lógica de negocio, la sincronización y la persistencia de la información.

Dentro del alcance del proyecto se incluyen las siguientes funcionalidades:

\begin{itemize}
    \item Diseño e implementación de la arquitectura completa del sistema,
    incluyendo aplicación móvil, portal web, backend y base de datos.

    \item Registro de observaciones de residuos utilizando una taxonomía
    jerárquica estandarizada que garantice la consistencia científica de los
    datos.

    \item Captura automática de información geográfica y temporal asociada a cada
    observación mediante los sensores disponibles en los dispositivos móviles.

    \item Gestión de fotografías y metadatos asociados a las observaciones
    registradas.

    \item Gestión de campañas y sesiones de muestreo, incluyendo información
    contextual como el entorno analizado, el método empleado, la duración de la
    sesión, el tipo de transecto o cuadrante y las características del área
    estudiada.

    \item Funcionamiento en modo offline mediante almacenamiento local de la
    información y sincronización posterior con el servidor cuando exista
    conectividad.

    \item Administración de campañas, usuarios, grupos y permisos mediante un
    portal web de gestión.

    \item Visualización de observaciones, sesiones y campañas mediante mapas
    interactivos y paneles informativos.

    \item Exportación de los datos recopilados en formatos adecuados para su
    posterior explotación y análisis científico.

    \item Implementación de mecanismos de autenticación, control de acceso y
    protección de la información almacenada.

    \item Elaboración de la documentación técnica, manuales, pruebas y memoria
    académica asociada al desarrollo del sistema.
\end{itemize}

Quedan fuera del alcance del Trabajo Fin de Grado el mantenimiento
post-producción, el soporte técnico continuo tras la entrega y la definición de
nuevos protocolos científicos oficiales. De esta forma, el alcance del proyecto
se centra en proporcionar una plataforma funcional para la recogida, gestión,
consulta y explotación de datos relacionados con la monitorización de residuos,
estableciendo una base tecnológica que permita futuras ampliaciones y desarrollos.